\subsection{The open-string \texorpdfstring{$KBc$}{KBc} dictionary}
\label{rs:sec:osft}

We now compare with Witten's open-string product and with the analytic
($KBc$) solution technology. The analogy is tight and entry-by-entry, and it does
useful work: it shows that manipulating a non-geometric object like $\mtwo^{-1}$
is exactly the closed-string counterpart of the wedge-algebra bookkeeping that
makes analytic solutions well-defined, and it locates the one genuinely open
question on the closed side.

\subparagraph{The $KBc$ algebra in one paragraph.} The analytic-solution technology
of open SFT \cite{Schnabl2005,ErlerSchnabl2009,Erler2019} is built on the algebra
$K,B,c$: $K$ \emph{generates} worldsheet strips, so $e^{-tK}$ is a wedge of width
$t$ (one surface per $t$, a continuous semigroup); $B$ is its $b$-ghost partner;
$c$ a boundary insertion. \Cref{rs:tab:osft} gives the entry-by-entry
correspondence.

\begin{table}[htp]
	\centering
	\renewcommand{\arraystretch}{1.3}
	\begin{tabular}{p{0.42\linewidth}p{0.42\linewidth}}
		\toprule
		open string ($KBc$) & closed HS \\
		\midrule
		$K$: strip \emph{generator}; $e^{-tK}$ = one surface per $t$ (continuous semigroup)
		& $L_0^+$: cylinder generator; $e^{-sL_0^+}$ = one surface per $s$ (both geometric) \\
		Witten $*$: half-string overlap, moduli-free vertex; covers $\mc{M}$ with $1/K$ alone \cite{GiddingsMartinecWitten1986}
		& vertex regions $\mc{V}_{0,n}$: moduli-free \emph{as vertices} (moduli pre-integrated); do \textbf{not} cover $\mc{M}$ alone (Sonoda--Zwiebach) --- hence the tower \\
		wedge algebra $e^{-t_1K}e^{-t_2K}=e^{-(t_1+t_2)K}$: width additive, continuous
		& chain monoid $\mtwo^{n}\mtwo^{n'}=\mtwo^{n+n'}$: length additive, \textbf{discrete} (\cref{rs:sec:frozen}) --- the central structural analogy \\
		$\dfrac{1}{1+K}=\int_0^\infty\!\dd t\,e^{-t}e^{-tK}$: geometric, every fibre an honest positive-width strip
		& $\mtwo^{-1}=\sum_n (D^{-1}N)^nD^{-1}$ \eqref{p:eq:ws:neumann}: ``geometric'', every term an honest frozen-neck chain \\
		inverse wedge $e^{+K}$ (``negative width''): formal, non-geometric; legal only with net positive width, $e^{K}e^{-2K}=e^{-K}$
		& $\mtwo^{-1}$ bare: formal, non-surface (\cref{p:sec:ws:inverse}); legal sandwiched, $V_4V_4^{-1}V_4=V_4$ \eqref{p:eq:hs4:telescoping} \\
		$1/K$ dangerous: $\spec(K)$ reaches $0$ (sliver); needs the security strip / phantom terms
		& $1/\mtwo$ conditional: $\spec(\mtwo)$ bounded from $0$ (\cref{ch:momentum-space}); kernel quotiented \\
		\bottomrule
	\end{tabular}
	\caption{The open-SFT $KBc$ dictionary. Left column: literature
	\cite{Schnabl2005,ErlerSchnabl2009,Erler2019,GiddingsMartinecWitten1986}.
	Right column: results of \cref{p:sec:hs:quartic,p:sec:ws:inverse,sec:ws1:schwinger}. The correspondence itself is an
	organising analogy, not a theorem.}
	\label{rs:tab:osft}
\end{table}

\subparagraph{Two uses of the dictionary.} \emph{(a)} It says exactly what an
analytic-solution practitioner means by manipulating non-geometric elements: what
the telescoping~\eqref{p:eq:hs4:telescoping} licenses --- formal inverses are permitted whenever the full expression
telescopes to net-geometric objects, and the amplitude identity \eqref{eq:ch2:amp:sum}
is the closed-string counterpart of the wedge-width bookkeeping that makes
solutions of the type $\psi=\frac{1}{1+K}c(1+K)Bc\frac{1}{1+K}$ well-defined.
\emph{(b)} It locates the genuinely open question: Erler-style control of
$1/(1+K)$ rests on $\spec(K)\ge0$; the closed analogue --- positivity or a gap of
$\mtwo$ on $\Sym^2\Hilb_0$ --- is the question of \cref{ch:momentum-space}, and the dictionary suggests
importing the OSFT toolkit (resolvent estimates, level-truncated spectral density)
rather than seeking a worldsheet representation of $\mtwo^{-1}$ that, by
\cref{p:sec:ws:inverse}, does not exist in the class of local kernels defined there. \emph{(Interpretive: the OSFT-side facts are
literature, the closed-side facts are the results cited in the table, the correspondence is an
analogy.)}

